\section{评估与 scaling：证据应该怎样读}

前面展示能力，本章检查证据。视频自动指标难以同时衡量真实感、运动自然、运动完整、文本遵循和美学，因此 Movie Gen 主要依赖人类成对比较；scaling 图则探索语言模型的 compute-optimal 趋势是否能近似预测媒体模型。两类证据都很重要，也都不能过度外推。

\subsection{Net Win Rate 如何计算}

前面看到的是精选样例，本节转向可聚合的人类证据。结果表将 Movie Gen 与 Runway Gen3、LumaLabs、OpenAI Sora 和 Kling1.5 比较。读表顺序应是：先看指标定义，再看正负号和置信区间，最后比较同一 baseline 下不同维度，不能只截取最大的百分比或把相对偏好误读成绝对质量。

\lecturefigure{slide-433-human-evaluation.jpg}{Movie Gen 对多种 prior work 的人类评测 Net Win Rate}{00:47:24}

对模型 A 与 B，令第 $i$ 个样本的共识得分 $s_i\in\{-1,0,1\}$，分别表示 B 胜、平局、A 胜，则
\[
\operatorname{NWR}(A,B)=\frac{100}{M}\sum_{i=1}^{M}s_i.
\]
$M$ 是评测样本数。正值表示 A 更常被偏好，负值表示 B 更常被偏好。论文通过 bootstrap 重采样 item-level scores 构造 95\% confidence interval。

\begin{warningbox}{Net Win Rate 不是绝对质量分}
它依赖 prompt 集、比较模型版本、评审标准和人群。A 对 B 为正、B 对 C 为正，并不保证 A 对 C 的差值可相加。表中不同维度还可能冲突，例如 motion completeness 改善但 aesthetics 下降。
\end{warningbox}

\teachervoice{00:45:13--00:47:24，老师说视频自动指标仍很弱，因此使用多维人类比较。他也提醒公平比较很难：模型接口、prompt rewrite、采样设置和可访问版本都会影响结论。}

\subsection{Movie Gen 与 Llama 3 scaling law 的关系}

下一张图不是“视频等于文本”的证明，而是一个经验惊喜：Movie Gen 在不同 compute budget 下的 optimal model size，与 Llama 3 拟合曲线相近。它支持统一 Transformer 可能共享某些尺度规律，但数据、token 语义和目标仍然不同。

\lecturefigure{slide-434-scaling-laws.jpg}{Movie Gen compute-optimal 模型规模与 Llama 3 scaling curve 的对照}{00:49:13}

可用幂律形式概括 compute-optimal 参数量：
\[
N_{\mathrm{opt}}(C)=aC^b,
\]
其中 $C$ 是训练 FLOPs，$N_{\mathrm{opt}}$ 是在该预算下最优参数量，$a,b$ 由实验拟合。图中蓝点来自 Movie Gen 小规模实验，橙线来自 Llama 3；相近不意味着 $a,b$ 在所有分辨率和数据分布都恒定。

\begin{importantbox}{Scaling law 的正确用途}
Scaling law 用于在预算内选择模型/数据配置和预测趋势，不用于替代最终高分辨率评测。它要求训练 recipe、数据质量和 loss 定义相对稳定；一旦 representation、curriculum 或数据过滤改变，旧拟合可能失效。
\end{importantbox}

\teachervoice{00:47:26--00:49:13，老师把 Llama 3 曲线称为 surprisingly reasonable predictor，但没有宣称媒体和文本完全共享定律。课程结论是“值得继续测”，不是“已经证明 modality independent”。}

\subsection{本章小结}

Movie Gen 用多维人类成对评测弥补自动指标不足，Net Win Rate 仍受 prompt、协议和 baseline 约束。Scaling-law 对照显示统一 Transformer 的经验趋势可能跨模态复用，但它只是特定区间的拟合，需要数据质量和训练 recipe 保持可比。

